% GENERATED FILE. Edit canonical lessons, then run npm run generate:books.
% canonical-source-hash: fnv1a64:3e096be6967db1c9
% canonical-lessons: TA-C298-vipattu, TA-C298-patukayamatai, TA-C298-urukkulai, TA-C298-uyirtappu, TA-C298-mutalutavi

\chapter{A Road Accident}
\label{ch:ta-a-road-accident}
\hlchaptermodality{298}

\begin{chapteropening}
\textbf{By the end of this chapter:} \emph{I can describe a road accident.}

Complete the last lesson of chapter 298: I can describe a road accident.
\end{chapteropening}

\section[\texorpdfstring{\ta{விபத்து} (vipattu)}{விபத்து (vipattu)}]{\ta{விபத்து} (vipattu) — an accident}

\label{lesson:TA-C298-vipattu}



\begin{quote}
Before the new one: say the Tamil for an investigation, then the Tamil for to accuse.
\end{quote}

\subsection*{You'll want to know: \ta{விபத்து}}
\textbf{\ta{விபத்து}} — \emph{vipattu} — \textquotedblleft{}an accident\textquotedblright{}.

\begin{quote}
\ta{வேகமாக} \ta{ஓட்டினால்} \ta{விபத்து} \ta{நடக்கும்}.
\end{quote}

\begin{quote}
\emph{vēkamāka ōṭṭiṉāl vipattu naṭakkum.}
\end{quote}

\begin{quote}
If you drive fast, an accident will happen.
\end{quote}

Plural \ta{விபத்துகள்}; a road accident is \ta{சாலை} \ta{விபத்து}.

\begin{cousinweb}
\textbf{\ta{விபத்து}} (\emph{vipattu}) is borrowed from Sanskrit \emph{vipatti}, \textquotedblleft{}misfortune, calamity\textquotedblright{} (CDIAL 11788); Tamil ends it in -\ta{உ}.
\end{cousinweb}

\begin{grammarlens}[title={What you've built}]
One more word for telling what happened.
\end{grammarlens}

\subsection*{Your turn}
\begin{itemize}
\raggedright
  \item \emph{Say it:} \emph{vipattu}
  \item \emph{Say it:} \emph{vipattu}, once more
  \item \emph{Say it:} \emph{vēkamāka ōṭṭiṉāl vipattu naṭakkum.} — If you drive fast, an accident will happen.
  \item \emph{Recall:} say the Tamil for an investigation, then the Tamil for to accuse, then say \emph{vipattu} again
\end{itemize}

\begin{tcolorbox}[breakable,colback=teal!4,colframe=teal!35!black,title={Before you move on}]
What does \ta{விபத்து} mean? (\textquotedblleft{}An accident\textquotedblright{}.) Say it once more.
\end{tcolorbox}
\section[\texorpdfstring{\ta{படுகாயமடை} (paṭukāyamaṭai)}{படுகாயமடை (paṭukāyamaṭai)}]{\ta{படுகாயமடை} (paṭukāyamaṭai) — to be badly hurt}

\label{lesson:TA-C298-patukayamatai}



\begin{quote}
Before the new one: say the Tamil for to accuse, then the Tamil for an accident.
\end{quote}

\subsection*{You'll want to know: \ta{படுகாயமடை}}
\textbf{\ta{படுகாயமடை}} — \emph{paṭukāyamaṭai} — \textquotedblleft{}to be badly hurt\textquotedblright{}.

\begin{quote}
\ta{விபத்தில்} \ta{ஓட்டுநர்} \ta{படுகாயமடைந்தார்}.
\end{quote}

\begin{quote}
\emph{vipattil ōṭṭunar paṭukāyamaṭaintār.}
\end{quote}

\begin{quote}
The driver was badly hurt in the accident.
\end{quote}

Past (I …): \ta{படுகாயமடைந்தேன்}, \emph{paṭukāyamaṭaintēṉ}.

\ta{படு}- before a noun makes it severe: \ta{படுகாயம்}, a severe wound.

\begin{cousinweb}
\textbf{\ta{படுகாயமடை}} (\emph{paṭukāyamaṭai}) is a Tamil compound built from \ta{படு}, \textquotedblleft{}severe\textquotedblright{}, \ta{காயம்}, \textquotedblleft{}a wound\textquotedblright{}, and \ta{அடை}, \textquotedblleft{}to reach, get\textquotedblright{} (DEDR 79).
\end{cousinweb}

\begin{grammarlens}[title={What you've built}]
One more word for telling what happened.
\end{grammarlens}

\subsection*{Your turn}
\begin{itemize}
\raggedright
  \item \emph{Say it:} \emph{paṭukāyamaṭai}
  \item \emph{Say it:} \emph{paṭukāyamaṭai}, once more
  \item \emph{Say it:} \emph{vipattil ōṭṭunar paṭukāyamaṭaintār.} — The driver was badly hurt in the accident.
  \item \emph{Recall:} say the Tamil for to accuse, then the Tamil for an accident, then say \emph{paṭukāyamaṭai} again
\end{itemize}

\begin{tcolorbox}[breakable,colback=teal!4,colframe=teal!35!black,title={Before you move on}]
What does \ta{படுகாயமடை} mean? (\textquotedblleft{}To be badly hurt\textquotedblright{}.) Say it once more.
\end{tcolorbox}
\section[\texorpdfstring{\ta{உருக்குலை} (urukkulai)}{உருக்குலை (urukkulai)}]{\ta{உருக்குலை} (urukkulai) — to be wrecked}

\label{lesson:TA-C298-urukkulai}



\begin{quote}
Before the new one: say the Tamil for an accident, then the Tamil for to be badly hurt.
\end{quote}

\subsection*{You'll want to know: \ta{உருக்குலை}}
\textbf{\ta{உருக்குலை}} — \emph{urukkulai} — \textquotedblleft{}to be wrecked\textquotedblright{}.

\begin{quote}
\ta{விபத்தில்} \ta{கார்} \ta{முழுவதும்} \ta{உருக்குலைந்தது}.
\end{quote}

\begin{quote}
\emph{vipattil kār muḻuvatum urukkulaintatu.}
\end{quote}

\begin{quote}
The car was completely wrecked in the accident.
\end{quote}

Past (it …): \ta{உருக்குலைந்தது}, \emph{urukkulaintatu}.

\ta{உரு}, shape + \ta{குலை}, to be upset: to lose its shape.

\begin{cousinweb}
\textbf{\ta{உருக்குலை}} (\emph{urukkulai}) is a Tamil compound built from \ta{உரு}, \textquotedblleft{}form, shape\textquotedblright{} (DEDR 657, which asks whether it comes from Sanskrit \emph{rūpa}), and \ta{குலை}, \textquotedblleft{}to be upset, thrown into disorder\textquotedblright{} (DEDR 1806).
\end{cousinweb}

\begin{grammarlens}[title={What you've built}]
One more word for telling what happened.
\end{grammarlens}

\subsection*{Your turn}
\begin{itemize}
\raggedright
  \item \emph{Say it:} \emph{urukkulai}
  \item \emph{Say it:} \emph{urukkulai}, once more
  \item \emph{Say it:} \emph{vipattil kār muḻuvatum urukkulaintatu.} — The car was completely wrecked in the accident.
  \item \emph{Recall:} say the Tamil for an accident, then the Tamil for to be badly hurt, then say \emph{urukkulai} again
\end{itemize}

\begin{tcolorbox}[breakable,colback=teal!4,colframe=teal!35!black,title={Before you move on}]
What does \ta{உருக்குலை} mean? (\textquotedblleft{}To be wrecked\textquotedblright{}.) Say it once more.
\end{tcolorbox}
\section[\texorpdfstring{\ta{உயிர்தப்பு} (uyirtappu)}{உயிர்தப்பு (uyirtappu)}]{\ta{உயிர்தப்பு} (uyirtappu) — to escape alive}

\label{lesson:TA-C298-uyirtappu}



\begin{quote}
Before the new one: say the Tamil for to be badly hurt, then the Tamil for to be wrecked.
\end{quote}

\subsection*{You'll want to know: \ta{உயிர்தப்பு}}
\textbf{\ta{உயிர்தப்பு}} — \emph{uyirtappu} — \textquotedblleft{}to escape alive\textquotedblright{}.

\begin{quote}
\ta{அதிர்ஷ்டவசமாக}, \ta{எல்லோரும்} \ta{உயிர்தப்பினார்கள்}.
\end{quote}

\begin{quote}
\emph{atirṣṭavacamāka, ellōrum uyirtappiṉārkaḷ.}
\end{quote}

\begin{quote}
Luckily, everyone escaped alive.
\end{quote}

Past (I …): \ta{உயிர்தப்பினேன்}, \emph{uyirtappiṉēṉ}.

Built from \ta{உயிர்}, life, and \ta{தப்பு}, to escape.

\begin{cousinweb}
\textbf{\ta{உயிர்தப்பு}} (\emph{uyirtappu}) is a Tamil compound built from \ta{உயிர்}, \textquotedblleft{}life\textquotedblright{} (DEDR 645), and \ta{தப்பு}, \textquotedblleft{}to escape, be saved\textquotedblright{} (DEDR 3071).
\end{cousinweb}

\begin{grammarlens}[title={What you've built}]
One more word for telling what happened.
\end{grammarlens}

\subsection*{Your turn}
\begin{itemize}
\raggedright
  \item \emph{Say it:} \emph{uyirtappu}
  \item \emph{Say it:} \emph{uyirtappu}, once more
  \item \emph{Say it:} \emph{atirṣṭavacamāka, ellōrum uyirtappiṉārkaḷ.} — Luckily, everyone escaped alive.
  \item \emph{Recall:} say the Tamil for to be badly hurt, then the Tamil for to be wrecked, then say \emph{uyirtappu} again
\end{itemize}

\begin{tcolorbox}[breakable,colback=teal!4,colframe=teal!35!black,title={Before you move on}]
What does \ta{உயிர்தப்பு} mean? (\textquotedblleft{}To escape alive\textquotedblright{}.) Say it once more.
\end{tcolorbox}
\section[\texorpdfstring{\ta{முதலுதவி} (mutalutavi)}{முதலுதவி (mutalutavi)}]{\ta{முதலுதவி} (mutalutavi) — first aid}

\label{lesson:TA-C298-mutalutavi}



\begin{quote}
Before the new one: say the Tamil for to be wrecked, then the Tamil for to escape alive.
\end{quote}

\subsection*{You'll want to know: \ta{முதலுதவி}}
\textbf{\ta{முதலுதவி}} — \emph{mutalutavi} — \textquotedblleft{}first aid\textquotedblright{}.

\begin{quote}
\ta{ஒரு} \ta{செவிலியர்} \ta{அவருக்கு} \ta{முதலுதவி} \ta{செய்தார்}.
\end{quote}

\begin{quote}
\emph{oru ceviliyar avarukku mutalutavi ceytār.}
\end{quote}

\begin{quote}
A nurse gave him first aid.
\end{quote}

The verb is \ta{செய்}: \ta{முதலுதவி} \ta{செய்}, to give first aid.

\begin{cousinweb}
\textbf{\ta{முதலுதவி}} (\emph{mutalutavi}) is a Tamil compound built from \ta{முதல்}, \textquotedblleft{}first\textquotedblright{} (DEDR 4950), and \ta{உதவி}, \textquotedblleft{}help\textquotedblright{} (DEDR 609).
\end{cousinweb}

\begin{grammarlens}[title={What you've built}]
One more word for telling what happened.
\end{grammarlens}

\subsection*{Your turn}
\begin{itemize}
\raggedright
  \item \emph{Say it:} \emph{mutalutavi}
  \item \emph{Say it:} \emph{mutalutavi}, once more
  \item \emph{Say it:} \emph{oru ceviliyar avarukku mutalutavi ceytār.} — A nurse gave him first aid.
  \item \emph{Recall:} say the Tamil for to be wrecked, then the Tamil for to escape alive, then say \emph{mutalutavi} again
\end{itemize}

\begin{tcolorbox}[breakable,colback=teal!4,colframe=teal!35!black,title={Before you move on}]
What does \ta{முதலுதவி} mean? (\textquotedblleft{}First aid\textquotedblright{}.) Say it once more.
\end{tcolorbox}
